\section{模型假设及代价}
\begin{enumerate}
\item 药材均质、轴对称，忽略轴向与周向梯度。长度与直径之比为6.25，两个端面与侧面面积之比为0.08，因而采用一维径向模型。端面附近的局部快速干燥不能由此准确刻画，本文时长应理解为该径向模型的预测。
\item 固体骨架与液态水局部热平衡，以单一温度场及题给有效物性描述内部导热；忽略辐射、化学反应和内部气相显式迁移。蒸发集中在表面，潜热作为边界能量通量计入，避免在内部再添加同一笔蒸发耗热。
\item 问题一至三尺寸不变，采用恒定参考干物质密度；问题四长度不变、径向均匀收缩，以附件2规定半径并保持总干物质量守恒。均匀收缩不能刻画内部裂隙、各向异性或局部塌陷，给定半径也不是模型自主预测的结果。
\item 经验密度$\rho$用于有效热容量，\rhod{}用于含水率与实际水分通量的换算。二者作用不同，不额外强迫题给经验密度同时满足每一时刻的总质量与体积关系；否则会与给定几何约束发生矛盾。
\item 附件数据在相邻时间点间分段线性插值，保存原始测量值且不产生额外极值。附件覆盖外的环境使用有文献依据的自主工况；所有外推、阶段切换和气膜闭合均在相应章节单独说明。
\item 未给出药材品种、品质劣变阈值、吸附实测曲线和风速场。所选工况与供水上限关系属于明确的模型闭合，预测时长具有条件性，不能据此声称对任何中药材均为品质或能耗最优。
\end{enumerate}
